\documentclass[11pt]{article}
\usepackage[margin=1in]{geometry}
\usepackage{amsmath,amssymb,amsthm}
\usepackage{booktabs}
\usepackage{graphicx}
\usepackage[colorlinks=true,linkcolor=blue,citecolor=blue,urlcolor=blue]{hyperref}
\usepackage{xcolor}

\newtheorem{lemma}{Lemma}
\newtheorem{proposition}{Proposition}
\newtheorem{definition}{Definition}

\title{Two Kinds of Saturation:\\ Why LLM-Judge Order Bias Is Hard to Measure}
\author{privatedick\thanks{Pseudonymous author. All code, proofs, and raw
results are public and reproducible at
\url{https://github.com/privatedick/llm-judge-saturation} (MIT); nothing in
this paper rests on trusting the author.}}
\date{August 2026}

\begin{document}
\maketitle

\begin{abstract}
Anyone who has used an LLM as a pairwise judge has watched it flip its
verdict when the two candidates are swapped. The standard remedy---run both
orders and compare---answers neither of the questions underneath: \emph{why}
does order matter, and why is \emph{measuring} the effect so easy to get
wrong? We separate two mechanisms that the word ``saturation'' conflates.
\emph{Internal} saturation (a clamped accumulator) can \textbf{create} order
sensitivity: we formalize a clamped belief update and prove, in Lean~4, that
constant-increment updates commute while the same update under a clamp does
not. \emph{Output} saturation (near-deterministic judge verdicts) does the
opposite---it \textbf{destroys} the observability that order statistics
require. Empirically, using a fully crossed $2\times2$ design that separates
content-position effects from label-preference effects (a confound in our own
first design, corrected here), we measure three open-weight judges and one
proprietary frontier judge: $64$--$68\%$ of judge conditions are saturated,
$14$ of $18$ (model, item) cells fail an exact-test interpretability gate,
effective sample sizes collapse to ${\approx}1.1$, and order effects are null
after Holm correction ($0/18$, replicated across two independent runs). On
the same tie items where order effects were near zero, label bias was
$7$--$10\times$ larger---item-specific, not a universal position prior. A
statistical power analysis shows the exact saturation gate has \emph{zero}
power below $k{=}29$ samples per condition at the tested threshold, so
quick-start sample sizes cannot in principle detect what they claim to
measure. The methodological point: before asking how large an order effect
is, ask whether the judge's output distribution contains any variance to
order. We release the crossed design, the gates, the Lean proofs (0 sorry,
axiom-audited at build time), and an offline proof-of-concept; total
experimental cost was \$0.46.
\end{abstract}

\section{Introduction}

LLM-as-judge evaluation asks a model ``which of these two answers is
better?'' and records the verdict. The verdict is not stable under swapping
the two candidates, and the field's standard fix is procedural: query both
orders, check agreement, average. The fix works well enough in practice, but
it leaves two questions unanswered. First, \emph{why} does order matter at
all---what is the mechanism? Second, why does \emph{measuring} the effect
keep going wrong in ways that the both-orders protocol does not catch?

This paper's claim is that the word ``saturation'' is doing double duty over
two mechanisms with opposite consequences:

\begin{enumerate}
\item \textbf{Internal saturation can manufacture an order effect.} A
bounded belief accumulator with a clamp---$x \mapsto \operatorname{clamp}(x
+ \delta)$---has state-dependent increments, and state-dependent increments
do not commute. We prove the positive and negative halves in Lean~4
(\S\ref{sec:lean}): clamped updates commute in the interior of the range and
provably fail to commute at the boundary.
\item \textbf{Output saturation can hide one---or make its measurement
meaningless.} A judge whose per-condition output distribution is
near-deterministic still produces a point verdict, but no variance, no error
bars, and effectively one observation. Order statistics computed on such a
distribution are non-identifying. Kang~\cite{kang2026} states this as a
caveat for QQ-equality audits; we make it operational as an exact test that
gates every downstream statistic (\S\ref{sec:design}).
\end{enumerate}

The empirical headline is a null result with a decomposition attached: after
separating position from label preference by full counterbalancing, we find
\emph{no} Holm-significant order effects in any of 18 cells across two
independent runs, on judges where $64$--$68\%$ of conditions are too
saturated for the question to be well-posed---while label bias on the same
items is an order of magnitude larger and item-specific. Our own first
design confounded the two signals and initially suggested an order effect;
the crossed rerun killed it. We report that correction in place
(\S\ref{sec:correction}) because the failure is the methodological point of
the paper.

\section{Background: order effects as non-commutation}

Write $T_A, T_B$ for the state updates of processing item $A$ then $B$. An
order effect is
\[
T_B(T_A(x)) \neq T_A(T_B(x)),
\]
i.e.\ the updates fail to commute. Quantum-cognition models instantiate this
with vectors and projectors and derive a parameter-free structural
invariant, the \textbf{QQ equality}~\cite{wang2013,wang2014}:
\[
p(A_y, B_n) + p(A_n, B_y) \;=\; p(B_y, A_n) + p(B_n, A_y),
\]
both orders agree on total ``changed my mind'' probability, however large
the individual order effect. Two caveats we carry throughout: satisfying QQ
does not establish the Hilbert-space model~\cite{boyerkassem2016}, and QQ is
one structural invariant among several, not a structure-vs-noise oracle.

\section{Internal saturation creates non-commutation}
\label{sec:lean}

Constant additive updates, $x \mapsto x + \delta$ for fixed $\delta$, always
commute---translations compose order-free. The cheapest way an
\emph{additive} update becomes order-sensitive is a clamp:
\begin{definition}[Clamped update]
$\operatorname{clamp}_{[lo,hi]}(y) = \max(lo, \min(hi, y))$, and the update
$U_\delta(x) = \operatorname{clamp}_{[lo,hi]}(x + \delta)$.
\end{definition}

Both halves are machine-checked in Lean 4 (repository,
\texttt{lean/EvoEcos/ClampedUpdateCommute.lean}, 0 sorry, with a build-time
axiom audit that fails the build on any axiom outside Lean's standard
three):

\begin{lemma}[Interior commutation]
If $x$, $x+a$, $x+a+b$ and their mirrored order all lie in the open interior
of $[lo,hi]$, then $U_b(U_a(x)) = U_a(U_b(x))$: within the interior the
clamped update is an affine translation and translations commute.
\end{lemma}

\begin{proposition}[Saturation breaks commutation]
There exist $lo, hi, x, a, b$ with the accumulator in range but at the
boundary where $U_b(U_a(x)) \neq U_a(U_b(x))$: the effective increment
becomes a function of state, and order sensitivity appears exactly where the
accumulator saturates.
\end{proposition}

The same repository proves the perturbation companion used for
framing/order sensitivity bounds (\texttt{PersuasionOperator.lean}): the
persuasion power satisfies $|\mathrm{power}(\varphi,\theta)| \le |\sin
\theta|$ and the bound is tight. The two results together say: internal
saturation is a \emph{sufficient} mechanism for order sensitivity (not a
claim that production LLMs implement clamped accumulators), and bounded
perturbation is the right frame for asking how far order or framing can move
a decision.

\section{Output saturation destroys observability}
\label{sec:saturation}

Kang~\cite{kang2026} audits question-order effects in LLMs with the QQ
equality and flags the caveat that a near-deterministic judge makes the QQ
comparison non-identifying: with no variance there is no
``changed-my-mind'' probability to compare across orders. We treat that
caveat as a first-class measurement gate rather than a footnote.

\begin{definition}[Saturation gate]
A condition's verdict distribution is \emph{saturated at level $s$} when the
top outcome probability exceeds $s$ under an exact one-sided binomial test
($H_0\colon p_{\text{top}} \le s$), not a bare point-estimate threshold: a
point-estimate check over-flags dispersed judges from finite-sample noise.
\end{definition}

We report saturation at two levels because they answer different questions:
the fraction of individual (position, label) conditions that are saturated,
and the fraction of (model, item) cells in which \emph{any} of the four
conditions is saturated---the conservative gate under which order statistics
for that cell are non-identifying.

The gate has a power floor that quick-start designs silently violate. At
$s = 0.9$ and $\alpha = 0.05$, the exact test has \emph{zero} power below
$k = 29$ samples per condition: at $k \le 28$ even an all-agree verdict
string (``AAAA\ldots'') cannot reject the null of non-saturation. The
crossover at $k=29$ is verified in the repository's test suite by running
the real exact test on synthetic all-agree lists at $k{=}28$ and $k{=}29$
(\texttt{python/test\_gate\_power.py}). Our first-pilot defaults sat below
the floor; the repository's earlier quick-start did too.

\section{Design: crossing position against label}
\label{sec:design}

A judge asked ``A or B?'' can prefer the first slot (a \emph{position}
effect) or prefer the token ``B'' (a \emph{label} effect). Every A/B-then-B/A
protocol without label counterbalancing measures their sum. The design used
here is a fully crossed $2\times2$ per item: (content order) $\times$
(label assignment), $k = 40$ samples per condition, temperature $1.0$.
Decomposition:

\begin{itemize}
\item \textbf{order effect} = the pure content-position signal (same labels,
swapped content), Fisher exact test per cell, Holm correction across cells;
\item \textbf{label bias} = the pure token-prior signal (same content,
swapped labels).
\end{itemize}

Three open-weight judges via OpenRouter (DeepSeek v4-flash, Qwen3-8b,
Mistral NeMo) on six items (three ``tie'' items with no clear winner, three
clear-winner controls), giving $18$ (model, item) cells and $72$ conditions
per run; two independent runs on different days. A proprietary frontier
judge (GPT-5-mini) was added later on a three-item subset. Total API cost
across all three runs: \$0.46.

\section{Results}

\begin{table}[h]
\centering
\small
\begin{tabular}{lcc}
\toprule
 & Run 1 (08-11) & Run 2 (08-12) \\
\midrule
conditions saturated ($s{=}0.9$, exact test) & $46/72$ ($63.9\%$) & $49/72$ ($68.1\%$) \\
cells failing interpretability gate & $14/18$ & $14/18$ \\
interpretable cells & $4$ & $4$ \\
mean order effect, all cells & $0.018$ & $0.024$ \\
mean order effect, interpretable & $0.041$ & $0.063$ \\
Holm-significant order effects & $0/18$ & $0/18$ \\
mean label bias, tie items & $0.456$ & $0.408$ \\
mean order effect, same tie items & $0.046$ & $0.058$ \\
mean effective sample size (worst condition) & $1.10$ & $1.10$ \\
estimated cost & \$0.21 & \$0.21 \\
\bottomrule
\end{tabular}
\caption{Two independent runs, three open-weight judges. The null replicates:
no cell shows an order effect surviving Holm correction; on the tie items,
label bias is $7$--$10\times$ the order effect measured on the same items.}
\end{table}

Three findings, in increasing order of interest:

\textbf{(1) The order-effect null replicates.} Zero of 18 cells per run,
twice. One cell in Run~2 reached raw $p < .05$ and did not survive Holm; at
this $k$, a difference of one sample is not a measured magnitude.

\textbf{(2) Most judge conditions are too saturated for the question to be
well-posed.} Between $64$ and $68\%$ of conditions and $14/18$ cells fail
the exact gate; effective sample size is ${\approx}1.1$. The judge returns a
verdict, but the verdict stream carries almost no information about
\emph{ordering} of anything.

\textbf{(3) Label bias dwarfs order bias and is item-specific.} On the tie
items---precisely where a position prior would express itself---mean label
bias was $0.41$--$0.46$ against order effects of $0.046$--$0.058$ on the same
items. The frontier judge sharpened the picture: GPT-5-mini showed label bias
of $0.49$ on one tie item (essentially unbiased) and $0.11$ on another
(strongly B-preferring)---the same judge, two items, opposite characters.
That is evidence for item-specific token quirks, not a universal position
prior.

\section{The correction, reported in place}
\label{sec:correction}

This paper's first design was wrong, and we report the failure because it is
the methodological point. The original A/B-then-B/A protocol confounded
content position with label preference; it initially suggested an order
effect. The crossed rerun decomposed the signal, and the effect moved
entirely into the label channel. Every number above is from the crossed
design.

The correction went both directions. An external review (2026-08-12) claimed
four reproductions of a parser bug and one missing length-check; when
actually run, three of the four parser examples returned the correct letter
against the code as shipped, and the length-check was already present and
already raising. The fixes that were real were applied only after
reproducing the failure directly---the same discipline this paper argues
for, applied to its own reviewers. Two review-requested extensions are also
reported honestly: the proprietary-judge extension succeeded (above); a
temperature sweep ($T \in \{0, 0.5, 1\}$) was attempted twice, failed twice
on the same infrastructure cause (a per-read inactivity timeout that a
trickling connection never trips; now a hard wall-clock deadline), and is
\emph{not} reported from incomplete runs. The temperature question stays
open.

\section{Limitations}

The judges are three open-weight models plus one frontier model on a
six-item panel: this is a measurement-methodology demonstration, not a
benchmark. The Lean results establish clamped saturation as a
\emph{sufficient} mechanism for non-commutation, not a claim about
transformer internals. $k=40$ per condition is above the gate's power floor
but modest for small effects. The QQ equality is used as a structural
reference point, and its satisfaction would not establish the Hilbert-space
model~\cite{boyerkassem2016}.

\section{Reproducibility}

Everything is public at
\url{https://github.com/privatedick/llm-judge-saturation} (MIT): the crossed
experiment and its tests, the exact gates with the $k{=}29$ power-floor
test, the Lean~4 proofs with the build-time axiom audit, raw result JSONs
for all runs, and an offline proof-of-concept requiring no API key. CI runs
the full suite on every push; the badges in the README are enforced, not
asserted. Total experimental cost: \$0.46.

\begin{thebibliography}{9}
\bibitem{wang2013} Z.~Wang and J.~R.~Busemeyer. A Quantum Question Order
Model Supported by Empirical Tests of Order Effects on Survey Measurements.
\emph{Topics in Cognitive Science} 5(4):689--710, 2013.
\bibitem{wang2014} Z.~Wang, T.~Solloway, R.~M.~Shiffrin, and J.~R.~Busemeyer.
Context effects produced by question orders reveal quantum nature of human
judgments. \emph{PNAS} 111(26):9431--9436, 2014.
\bibitem{boyerkassem2016} M.~Boyer-Kassem, S.~Duch\^ene, and E.~Guerci.
Testing quantum-like models of judgment for question order effect.
arXiv:1501.04901, 2016.
\bibitem{kang2026} P.~Kang. Auditing Question-Order Effects in Large
Language Models with the QQ Equality: Mechanism Characterization and a
Saturation Caveat. arXiv:2607.17219, 2026.
\end{thebibliography}

\end{document}
